\section{The main equation: twenty-three unrestricted integers}
\label{sec:integer}

\subsection{Fixed numerals and the complete polynomial}

Use the coherent source-34 tuple, named
\texttt{mixed-parity-unrounded-1064921}, together with source 71's two
determinant scales:
\[
 \beta=2^{1064921}-1,\qquad A_0=2^{1113},\qquad B_0=2^{666104}.
\]
The remaining three native numerals are $K_0,M_C,M_F$. Their complete
integer bytes, their hashes and their native binding are in
\texttt{integer/coefficients/BINDING.json} and its six payloads.
$M_F$ is the already shifted source mask. These symbols are fixed
coefficients, not unknowns. They are the exact integers of the pinned
source, not arbitrary integers satisfying some inequalities.

The twenty-three unknowns, all unrestricted integers, are
\begin{equation}
 J,F,\alpha,f,h,r,v,\tau,\eta,y,\gamma,o,j,E,c,a,I,u,Z,t,p_1,p_2,p_3.
 \label{eq:vars23}
\end{equation}
In the data files, $v,t,p_1,p_2,p_3$ are named
\texttt{v\_scale}, \texttt{pair\_root}, and
\texttt{pair\_pos1}, \texttt{pair\_pos2}, \texttt{pair\_pos3}.
Define only polynomial abbreviations:
\begin{align*}
 q&=\beta J+1,& C&=q-F-\alpha,&Y&=a-E,\\
 H&=4a+3,&\Delta&=a^2+H,&X&=4qu^2,\\
 D&=X+ac+\gamma H,&V&=fo-c,&k&=1-V+jc+hE,\\
 \zeta&=k-\eta,&\Pi&=1+p_1^2+p_2^2+p_3(p_3+1).
\end{align*}
Also put
\begin{align*}
 R&=\bigl((q-1)(q-F)+\alpha\bigr)(q^2-1)+(M_C+qM_F)J,\\
 G&=(A_0F)(B_0C)-(q-A_0F-B_0C)^2.
\end{align*}
The eleven residuals are
\begin{align}
 r_1&=Y-(2v+1)q^3,\nonumber\\
 r_2&=E-XY,\nonumber\\
 r_3&=c-\eta-kY,\nonumber\\
 r_4&=\tau^2+(2\tau-k)E(c-\eta)-1,\nonumber\\
 r_5&=D^2-\Delta c^2-1,\nonumber\\
 r_6&=jc-V-R,\nonumber\\
 r_7&=(K_0q+X)C-qF-Z(q-1),\nonumber\\
 r_8&=I-(2r+1)c^2,\nonumber\\
 r_9&=I^2-\Delta(f^2-1),\nonumber\\
 r_{10}&=y^2+I^2(V^2-y^2)-1,\nonumber\\
 r_{11}&=G\eta\zeta-(t^2-G-\eta\zeta+1)^2-\Pi.
 \label{eq:23residuals}
\end{align}
The proposed single Diophantine equation is
\begin{equation}
 \boxed{\displaystyle P_{23}=\sum_{\nu=1}^{11}r_\nu^2=0,
 \qquad\mathbf x\in\Z^{23}.}
 \label{eq:P23}
\end{equation}
All abbreviations are to be substituted; none introduces another
existential coordinate. Every exponent involving an unknown is fixed.
The companion JSON circuit evaluates precisely this polynomial in
122 live binary arithmetic operations.

\begin{theorem}[The 23-integer reduction]
The polynomial \eqref{eq:P23} has exact total degree twelve.
Its integer solvability is equivalent to the existence of a finite ZFC
refutation through the pinned native/compiler calibration of
sources 64, 71 and 75. Consequently its unsolvability has the consistency
interpretation \eqref{eq:target}, with the proof boundary stated in
Section 1.
\end{theorem}

We prove the new parts in detail, and specify the inherited interface at
the point where it is used.

\subsection{One certificate for two positive quantities}

For integers $b_1,b_2,b_3$ put
\[
 \Pos(b_1,b_2,b_3)=1+b_1^2+b_2^2+b_3(b_3+1).
\]
This is always strictly positive. Conversely every positive integer $n$
has this form: $4n-3$ is positive and congruent to $1$ or $5$ modulo $8$,
so Legendre's three-square theorem applies. Modulo $4$, exactly one of
the three squares is odd. Writing its root as $2b_3+1$ and the other two
as $2b_1,2b_2$ proves the assertion.

\begin{lemma}[Shared positivity, source 68]
\label{lem:shared}
For integers $x,y$, one has $x>0$ and $y>0$ if and only if there are
integers $t,b_1,b_2,b_3$ such that
\begin{equation}
 xy-(t^2-x-y+1)^2=\Pos(b_1,b_2,b_3).
 \label{eq:pair}
\end{equation}
\end{lemma}
\begin{proof}
If the left side is positive, then $xy>0$. Both numbers therefore have
the same nonzero sign. They cannot both be negative: writing
$x=-a,y=-b$, with $a,b\ge1$, makes the subtracted square at least
$(a+b+1)^2>ab$.

Conversely suppose $x,y\ge1$ and set $N=x+y-1$. Then
$xy-N=(x-1)(y-1)\ge0$. Let $m=\lfloor\sqrt N\rfloor$. Choose $t=m$
if $N-m^2\le m$, and $t=m+1$ otherwise. If $N$ is a square the error
is zero. Otherwise $|t^2-N|\le m$ and $m^2<N$, so
$(t^2-N)^2<N\le xy$. The resulting positive difference has a
$\Pos$ representation. The square root and floor select witnesses in
the proof; they are not operations in the polynomial.
\end{proof}

Apply this to $x=G$, $y=\eta\zeta$. The final row in
\eqref{eq:23residuals} therefore enforces exactly
\[
 G>0,\qquad\eta\zeta>0.
\]
It uses four private coordinates instead of the six needed by two
independent three-square representations. The equivalence holds for every
integer assignment, including assignments with a negative radix.

\subsection{The determinant and its canonical calibration}

Set $U=A_0F$ and $W=B_0C$. If $G=UW-(q-U-W)^2>0$, then $U,W$ have
one common nonzero sign and
\[
 |q-U-W|<\sqrt{UW}\le\frac{|U|+|W|}{2}.
\]
It follows that $q$ has that common sign and
\[
 |q|>\frac{A_0|F|+B_0|C|}{2}\ge |F|+|C|.
\]
Thus $F,C$ and $\alpha=q-F-C$ all have the same strict sign as $q$.
In particular $F\alpha>0$ and $FC>0$. A new zero can therefore restore
the two older native sign guards by independent $\Pos$ representations.
This implication is elementary and uses no native compiler theorem.

The converse uses canonical witnesses; the determinant is not intended
to contain every positive source assignment. Source 71 proves that the
canonical native accepting words for this exact tuple satisfy
\begin{equation}
 \frac q2\le A_0F<\frac{4q}{5},\qquad
 \frac q2\le B_0C<\frac{3q}{5},\qquad q>0.
 \label{eq:canonicalcone}
\end{equation}
Here is the numerical content of that calibration, so its dependence is
auditable. The native radix is $\rho=2^{121}$ and the cell length is
8801. Every occupied canonical cell has top $C$ digit $1$ at position
3296. Its top $F$ digit at position 8791 is
\[
 c_0=105534389824824194647382298625,
 \qquad 2^{96}\le c_0<c_0+1\le3\cdot2^{95}.
\]
The carry-free block bounds are
\begin{align*}
 \rho^{3296}&\le C_j<(1+1/\rho)\rho^{3296},\\
 c_0\rho^{8791}&\le F_j<(c_0+1)\rho^{8791}.
\end{align*}
The exponent identities
\[
 121\cdot8791+1113=1064921-97,\qquad
 121\cdot3296+666104=1064921-1
\]
give \eqref{eq:canonicalcone} after summing the canonical cells.
The full digit-support and carry proof is retained as
\texttt{integer/upstream/DETERMINANT\_GUARD.md}.

To see that these bounds imply $G>0$, write $s=q/2$,
$A_0F=s+x$, $B_0C=s+y$. Then $0\le x<3s/5$ and
$0\le y<s/5$, whence
\[
 G=s^2+s(x+y)-(x^2+y^2+xy)>\frac{12}{25}s^2>0.
\]
Source 71's bounds explicitly allow the lower dummy digits used to tune
the rank. Source 75's signed-bootstrap change modifies only its square
witness and cleared transport quotient, leaving $J,F,C,\alpha$
unchanged. Thus the canonical determinant and the signed bootstrap apply
to the same family of witnesses. This compatibility is needed for
completeness and was checked separately.

\subsection{Restoring the signed source}

Source 75 supplies a forty-integer-variable source with the arithmetic
constraints $X=4qu^2$ and
\[
 (K_0q+X)C-qF-Z(q-1)=0.
\]
Its signed Pell argument recovers $X=2^R>0$, $q>0$ and the odd rank
$R$, without assuming these signs at the start. The transport equation
then gives $q\mid Z$, restoring the older quotient. Canonical
completeness has $q=2^s$, with $s$ and $R$ odd and $R-s-2\ge0$,
so $u=2^{(R-s-2)/2}$ exists. This inherited existential source reduction
is retained in \texttt{integer/upstream/SIGNED\_BOOTSTRAP.md}.

To restore its literal arithmetic from a zero of $P_{23}$, use
\begin{align*}
 L&=q^2-1,&N&=q^3,&P&=k+E(c-\eta),\\
 C_2&=c^2,&F_2&=f^2,&g&=y^2-V^2,&B_{\rm aux}&=I^2,
\end{align*}
together with the abbreviations already displayed and the root shear
\[
 \tau_{\rm old}=\tau+E(c-\eta).
\]
All these restored coordinates are integer polynomials. The shear is
exact because, writing $P-k=E(c-\eta)$,
\[
 (\tau+P-k)^2-P^2+Pk-1
 =\tau^2+(2\tau-k)E(c-\eta)-1.
\]
Twelve of the signed source's twenty-two arithmetic rows become zero
polynomial identities; the other ten become $r_1,\ldots,r_{10}$ up to
sign. The independent audit verifies this entire identity list.

The last row restores $G>0$ and $\eta\zeta>0$ by
Lemma~\ref{lem:shared}. The determinant then restores the old
$F\alpha$ and $FC$ sign guards. This gives a signed-source zero and
therefore a finite refutation by its inherited soundness theorem.

Conversely, a finite refutation supplies canonical signed-source
witnesses. The determinant is positive on them by
\eqref{eq:canonicalcone}, and $\eta\zeta>0$. Select the four shared
positivity witnesses, apply the inverse root shear, and project away the
monically defined coordinates. Every residual in \eqref{eq:23residuals}
vanishes. This proves the two existence implications. The determinant
replacement is an existence equivalence using canonical witnesses;
the subsequent monic eliminations themselves are bijections of zero sets.

\subsection{Exact degree, not a gate-propagation estimate}

In the displayed order the residual degrees are
\[
 4,4,3,4,6,4,4,3,4,6,6.
\]
Their sum of squares has the following complete degree-twelve homogeneous
part, where every supplied unknown has degree one:
\begin{equation}
 \boxed{\lead_{12}(P_{23})=
 256\beta^4J^4u^8+I^4f^4o^4
       +\eta^4(-fo+jc+hE)^4.}
 \label{eq:leading23}
\end{equation}
Indeed the degree-six parts of $r_5,r_{10},r_{11}$ are respectively
$16\beta^2J^2u^4$, $I^2f^2o^2$, and
$-\eta^2(-fo+jc+hE)^2$. Every other residual has lower degree.
In particular
\[
 [\eta^4j^4c^4]P_{23}=1.
\]
This proves exact degree twelve after the authentic fixed integers are
substituted. It cannot be invalidated by a specialization relation among
the coefficient ports.

\subsection{The sextic and octic companions}

The elimination chain also yields useful lower-degree alternatives.
Starting from source 75's centered 34-variable sextic, replace its two
native sign guards by the determinant guard. This deletes six private
unknowns and introduces three, giving a 31-variable, 18-residual sextic.
Its coefficient $[r^2c^4]=4$ certifies the degree. Its supplied circuit
has 144 operations.

Eliminate $L=q^2-1$, $\Delta=a^2+4a+3$, and $B_{\rm aux}=I^2$.
Replace the two remaining positivity representations by the shared
certificate while retaining $\zeta$ as a coordinate. This gives
$31-3-6+4=26$ integer unknowns, fourteen residuals, degree eight and
131 operations. The private monomial $t^8$ has coefficient $1$.

Finally eliminate
\[
 \zeta=1-V+jc+hE-\eta,\quad
 D=4qu^2+ac+\gamma(4a+3),\quad V=fo-c.
\]
The first elimination makes the shared guard sextic. The other two make
the main and auxiliary norms sextic as well, so they cost no further
increase in the final degree. These are the three eliminations producing
\eqref{eq:P23}. The retained 25-variable intermediate documents this
chain, but is dominated by the 23-variable endpoint and is not counted
as a separate frontier result.
